\section{Literature review and research position}

\subsection{Risk-aware crop planning}

Crop-planning models have long represented planting as an operational
allocation rather than a list of recommended crops. Their variables include
crop areas, sequences and timing; their constraints include land, capital,
labour, equipment, rotations, storage and contracts. Contemporary models
combine these objects with crop-selection integer programmes
\citep{filippi2017mixed}, dynamic rotation benefits
\citep{boyabatli2019crop}, sustainable rotation systems
\citep{benini2023crop}, contract-farming schedules \citep{li2015contract} and
spatial decision support \citep{pahmeyer2021fruchtfolge}. Reviews show that
climate uncertainty is handled through stochastic, robust and multiobjective
approaches rather than one standardized decision model
\citep{randall2022robust,randall2024systematic}.

Loss-CVaR is useful in this setting because it controls expected shortfall and
admits a linear finite-scenario representation
\citep{rockafellar2000optimization,rockafellar2002conditional}. Agricultural
precedent is substantial but fragmented. \citet{aitsahlia2011enso} combine
ENSO-conditioned planting and hedging with Gaussian-copula scenarios and
CVaR. \citet{filippi2017mixed} embed CVaR in operational crop selection.
\citet{larsen2015geographical} and \citet{nguyenhuy2018copula} use copulas and
mean-CVaR to study geographical wheat portfolios. These studies establish
that crop portfolios can be optimized under downside risk. The present study
extends this perspective by comparing an externally supplied performance
ranking with optimized acreage, characterizing reversal boundaries and
examining the complete objective-equivalent optimal face.

\subsection{Joint tail risk and the limits of conventional diversification}

Diversification is effective when crop outcomes do not fail together in the
states that matter to the decision. Linear correlation summarizes average
co-movement, not the geometry of the joint lower tail. Student-\(t\) copulas
retain nonzero symmetric tail dependence, while Archimedean families can
concentrate dependence asymmetrically \citep{demarta2005tcopula}. Dependence
comparisons are family- and order-specific \citep{ansari2024dependence}; a
single lower-tail coefficient does not order arbitrary portfolio losses
across families.

Agricultural evidence motivates this distinction. Joint yield tails vary with
land quality \citep{du2018land}, production shocks are dependent across major
breadbaskets \citep{pflug2017breadbaskets}, and copula estimates show that the
risk reduction from crop diversification changes with crop pair, region and
drought severity \citep{schmitt2024drought}. These findings do not imply
general diversification failure.

Crop count and variance reduction are incomplete diagnostics for a
loss-CVaR-constrained portfolio. The diversification criterion used here
requires a conventional policy to pass a fixed variance benchmark yet fail
under the Student-\(t\) evaluation law and differ from the tail-aware
allocation.

\subsection{Prediction, prescription and performance scores}

Agricultural machine learning commonly predicts yield, classifies land or
ranks crop performance using weather, soil, management and remote-sensing
features \citep{liakos2018machine,vanklompenburg2020crop}. A score can be
scientifically useful even when it is not an acreage prescription. The
prescriptive-analytics literature instead evaluates how conditional
predictions change an optimized decision
\citep{bertsimas2020predictive,elmachtoub2022smart}. Decision-focused learning
extends this principle to end-to-end systems
\citep{mandi2023decision}.

The present analysis does not develop a learning algorithm. Instead, it
separates prediction from prescription. A pre-decision historical
relative-yield performance index ranks crop outcomes but does not enter the
profit objective. Optimized acreage can therefore be compared with the
external ranking without treating predictive accuracy as evidence of
prescriptive optimality or requiring the index to contain every economic
variable.

\subsection{Information value and operational flexibility}

Seasonal climate information has value only through decisions it can change.
Reviews emphasize that forecast skill, risk preferences, institutions and
available management responses jointly determine economic value
\citep{meza2008economic}. Bioeconomic applications measure value through
changes in sowing and other decisions \citep{anvo2021framework}, while
ENSO-based crop-planning models combine forecasts with planting and financial
hedges \citep{aitsahlia2011enso}. General decision analysis likewise links
information value to the portfolio or action set
\citep{keisler2004value,merkhofer1975value}; operational-flexibility theory
shows that flexibility value depends on margins, cost and the full uncertainty
structure \citep{vanmieghem1998investment}.

These results do not imply unconditional strict complementarity. A flexible action
can make a signal actionable, but it can also provide a robust fallback that
reduces the need to distinguish states. We therefore prove non-negativity,
strictness and zero conditions separately, add an increasing-differences
condition for complementarity, and construct an agricultural buffering path
for substitution.

\subsection{Contribution matrix}

\begin{table}[t]
\centering\scriptsize
\caption{\textbf{Positioning relative to selected research streams.}
\(\bullet\) denotes a central object, \(\circ\) denotes conditional or
subfield-specific coverage, and \(--\) denotes an object outside the stream's
principal scope. The comparison is selective rather than an exhaustive map of
the literature.}
\begin{tabularx}{\textwidth}{>{\raggedright\arraybackslash}Xcccccc}
\toprule
Stream & Performance rank & Operations & Loss-CVaR & Joint tail law &
Reversal frontier & Info.--flex.\\
\midrule
Agricultural ML reviews & \(\bullet\) & \(--\) & \(--\) & \(--\) & \(--\) & \(--\)\\
Crop-planning optimization & \(--\) & \(\bullet\) & \(\circ\) & \(\circ\) & \(--\) & \(--\)\\
Agricultural copula--CVaR & \(--\) & \(\circ\) & \(\bullet\) & \(\bullet\) & \(--\) & \(--\)\\
Climate forecast valuation & \(--\) & \(\circ\) & \(\circ\) & \(--\) & \(--\) & \(\circ\)\\
Present study & \(\bullet\) & \(\bullet\) & \(\bullet\) &
\(\bullet\) & \(\bullet\) & \(\bullet\)\\
\bottomrule
\end{tabularx}
\end{table}

Table~1 identifies how the study combines established research strands; it
does not imply exhaustive novelty. The model connects a historical
relative-yield performance index to stochastic operational allocation,
characterizes a reversal frontier, distinguishes variance and downside-risk
criteria for diversification, and identifies conditions under which
information complements or substitutes for flexibility. The next section
formalizes these elements in a common crop-planning model.
